% GENERATED FILE. Edit canonical lessons, then run npm run generate:books.
% canonical-source-hash: fnv1a64:95a545b37bbfe0db
% canonical-lessons: UR-C466-bazu, UR-C466-ran, UR-C466-pindli, UR-C466-talva, UR-C466-panja

\chapter{Arms and Legs}
\label{ch:ur-arms-and-legs}
\hlchaptermodality{466}

\begin{chapteropening}
\textbf{By the end of this chapter:} \emph{I can name more parts of the arms and legs.}

Complete the last lesson of chapter 466: I can name more parts of the arms and legs.
\end{chapteropening}

\section[\texorpdfstring{\ur{بازو} (bāzū)}{بازو (bāzū)}]{\ur{بازو} (bāzū) — an arm}

\label{lesson:UR-C466-bazu}



\begin{quote}
Before the new one: say the Urdu for the backbone, then the Urdu for a jaw.
\end{quote}

\subsection*{You'll want to know: \ur{بازو}}
\textbf{\ur{بازو}} — \emph{bāzū} — \textquotedblleft{}an arm\textquotedblright{}.

\begin{quote}
\ur{بھائی} \ur{کے} \ur{بازو} \ur{بڑے} \ur{مضبوط} \ur{ہیں۔}
\end{quote}

\begin{quote}
\emph{bhāī ke bāzū baṛe mazbūt haiṅ.}
\end{quote}

\begin{quote}
My brother's arms are very strong.
\end{quote}

Masculine.

\begin{cousinweb}
\textbf{\ur{بازو}} (\emph{bāzū}) is usually traced to Persian \emph{bāzū}, \textquotedblleft{}arm\textquotedblright{}.
\end{cousinweb}

\begin{grammarlens}[title={What you've built}]
One more word for describing what you lived and felt.
\end{grammarlens}

\subsection*{Your turn}
\begin{itemize}
\raggedright
  \item \emph{Say it:} \emph{bāzū}
  \item \emph{Say it:} \emph{bāzū}, once more
  \item \emph{Say it:} \emph{bhāī ke bāzū baṛe mazbūt haiṅ.} — My brother's arms are very strong.
  \item \emph{Recall:} say the Urdu for the backbone, then the Urdu for a jaw, then say \emph{bāzū} again
\end{itemize}

\begin{tcolorbox}[breakable,colback=teal!4,colframe=teal!35!black,title={Before you move on}]
What does \ur{بازو} mean? (\textquotedblleft{}An arm\textquotedblright{}.) Say it once more.
\end{tcolorbox}
\section[\texorpdfstring{\ur{ران} (rān)}{ران (rān)}]{\ur{ران} (rān) — a thigh}

\label{lesson:UR-C466-ran}



\begin{quote}
Before the new one: say the Urdu for a jaw, then the Urdu for an arm.
\end{quote}

\subsection*{You'll want to know: \ur{ران}}
\textbf{\ur{ران}} — \emph{rān} — \textquotedblleft{}a thigh\textquotedblright{}.

\begin{quote}
\ur{گرنے} \ur{سے} \ur{اس} \ur{کی} \ur{ران} \ur{پر} \ur{چوٹ} \ur{آئی۔}
\end{quote}

\begin{quote}
\emph{girne se us kī rān par choṭ āī.}
\end{quote}

\begin{quote}
He hurt his thigh when he fell.
\end{quote}

Feminine.

\begin{cousinweb}
\textbf{\ur{ران}} (\emph{rān}) is usually traced to Persian \emph{rān}, \textquotedblleft{}thigh\textquotedblright{}.
\end{cousinweb}

\begin{grammarlens}[title={What you've built}]
One more word for describing what you lived and felt.
\end{grammarlens}

\subsection*{Your turn}
\begin{itemize}
\raggedright
  \item \emph{Say it:} \emph{rān}
  \item \emph{Say it:} \emph{rān}, once more
  \item \emph{Say it:} \emph{girne se us kī rān par choṭ āī.} — He hurt his thigh when he fell.
  \item \emph{Recall:} say the Urdu for a jaw, then the Urdu for an arm, then say \emph{rān} again
\end{itemize}

\begin{tcolorbox}[breakable,colback=teal!4,colframe=teal!35!black,title={Before you move on}]
What does \ur{ران} mean? (\textquotedblleft{}A thigh\textquotedblright{}.) Say it once more.
\end{tcolorbox}
\section[\texorpdfstring{\ur{پنڈلی} (pinḍlī)}{پنڈلی (pinḍlī)}]{\ur{پنڈلی} (pinḍlī) — the calf of the leg}

\label{lesson:UR-C466-pindli}



\begin{quote}
Before the new one: say the Urdu for an arm, then the Urdu for a thigh.
\end{quote}

\subsection*{You'll want to know: \ur{پنڈلی}}
\textbf{\ur{پنڈلی}} — \emph{pinḍlī} — \textquotedblleft{}the calf of the leg\textquotedblright{}.

\begin{quote}
\ur{دوڑنے} \ur{کے} \ur{بعد} \ur{میری} \ur{پنڈلی} \ur{میں} \ur{درد} \ur{ہے۔}
\end{quote}

\begin{quote}
\emph{dauṛne ke bād merī pinḍlī meṅ dard hai.}
\end{quote}

\begin{quote}
My calf hurts after running.
\end{quote}

Feminine.

\begin{cousinweb}
\textbf{\ur{پنڈلی}} (\emph{pinḍlī}) is inherited: CDIAL lists the Hindi cousin \emph{piṇḍlī}, \textquotedblleft{}calf of the leg\textquotedblright{}, under Sanskrit \emph{píṇḍa}, \textquotedblleft{}a lump\textquotedblright{}, which already meant the calf (CDIAL 8168).
\end{cousinweb}

\begin{grammarlens}[title={What you've built}]
One more word for describing what you lived and felt.
\end{grammarlens}

\subsection*{Your turn}
\begin{itemize}
\raggedright
  \item \emph{Say it:} \emph{pinḍlī}
  \item \emph{Say it:} \emph{pinḍlī}, once more
  \item \emph{Say it:} \emph{dauṛne ke bād merī pinḍlī meṅ dard hai.} — My calf hurts after running.
  \item \emph{Recall:} say the Urdu for an arm, then the Urdu for a thigh, then say \emph{pinḍlī} again
\end{itemize}

\begin{tcolorbox}[breakable,colback=teal!4,colframe=teal!35!black,title={Before you move on}]
What does \ur{پنڈلی} mean? (\textquotedblleft{}The calf of the leg\textquotedblright{}.) Say it once more.
\end{tcolorbox}
\section[\texorpdfstring{\ur{تلوا} (talvā)}{تلوا (talvā)}]{\ur{تلوا} (talvā) — the sole of the foot}

\label{lesson:UR-C466-talva}



\begin{quote}
Before the new one: say the Urdu for a thigh, then the Urdu for the calf of the leg.
\end{quote}

\subsection*{You'll want to know: \ur{تلوا}}
\textbf{\ur{تلوا}} — \emph{talvā} — \textquotedblleft{}the sole of the foot\textquotedblright{}.

\begin{quote}
\ur{گرم} \ur{فرش} \ur{پر} \ur{اس} \ur{کے} \ur{تلوے} \ur{جل} \ur{گئے۔}
\end{quote}

\begin{quote}
\emph{garm farsh par us ke talve jal gae.}
\end{quote}

\begin{quote}
The hot floor burnt the soles of his feet.
\end{quote}

Masculine.

\begin{cousinweb}
\textbf{\ur{تلوا}} (\emph{talvā}) is inherited from \emph{talapāda}, \textquotedblleft{}the lower foot\textquotedblright{}, a form CDIAL reconstructs (CDIAL 5738).
\end{cousinweb}

\begin{grammarlens}[title={What you've built}]
One more word for describing what you lived and felt.
\end{grammarlens}

\subsection*{Your turn}
\begin{itemize}
\raggedright
  \item \emph{Say it:} \emph{talvā}
  \item \emph{Say it:} \emph{talvā}, once more
  \item \emph{Say it:} \emph{garm farsh par us ke talve jal gae.} — The hot floor burnt the soles of his feet.
  \item \emph{Recall:} say the Urdu for a thigh, then the Urdu for the calf of the leg, then say \emph{talvā} again
\end{itemize}

\begin{tcolorbox}[breakable,colback=teal!4,colframe=teal!35!black,title={Before you move on}]
What does \ur{تلوا} mean? (\textquotedblleft{}The sole of the foot\textquotedblright{}.) Say it once more.
\end{tcolorbox}
\section[\texorpdfstring{\ur{پنجہ} (panjā)}{پنجہ (panjā)}]{\ur{پنجہ} (panjā) — a paw; an outspread hand}

\label{lesson:UR-C466-panja}



\begin{quote}
Before the new one: say the Urdu for the calf of the leg, then the Urdu for the sole of the foot.
\end{quote}

\subsection*{You'll want to know: \ur{پنجہ}}
\textbf{\ur{پنجہ}} — \emph{panjā} — \textquotedblleft{}a paw; an outspread hand\textquotedblright{}.

\begin{quote}
\ur{بلی} \ur{نے} \ur{پنجے} \ur{سے} \ur{چوہا} \ur{پکڑا۔}
\end{quote}

\begin{quote}
\emph{billī ne panje se chūhā pakṛā.}
\end{quote}

\begin{quote}
The cat caught the mouse with its paw.
\end{quote}

Masculine.

\begin{cousinweb}
\textbf{\ur{پنجہ}} (\emph{panjā}) is usually traced to Persian \emph{panja}, \textquotedblleft{}the hand with its five fingers\textquotedblright{}, built on \emph{panj}, \textquotedblleft{}five\textquotedblright{}.
\end{cousinweb}

\begin{grammarlens}[title={What you've built}]
One more word for describing what you lived and felt.
\end{grammarlens}

\subsection*{Your turn}
\begin{itemize}
\raggedright
  \item \emph{Say it:} \emph{panjā}
  \item \emph{Say it:} \emph{panjā}, once more
  \item \emph{Say it:} \emph{billī ne panje se chūhā pakṛā.} — The cat caught the mouse with its paw.
  \item \emph{Recall:} say the Urdu for the calf of the leg, then the Urdu for the sole of the foot, then say \emph{panjā} again
\end{itemize}

\begin{tcolorbox}[breakable,colback=teal!4,colframe=teal!35!black,title={Before you move on}]
What does \ur{پنجہ} mean? (\textquotedblleft{}A paw; an outspread hand\textquotedblright{}.) Say it once more.
\end{tcolorbox}
